\section{Proofs for selected SOS formats and conservative extensions}

\begin{annexitem}{annex:p09-sos-local-semantics}{thm:p09-sos-local-semantics}{The source-size semantic certificate, with every negative target}
\begin{proof}
Rank every action or predicate judgement by the finite node count
of its source. In a ground local rule its source is \(f(p_1,\ldots,p_n)\).
Every positive-premise source and every possible denier of every
negative-premise source is one of the \(p_i\), a proper subterm.
Hence all required inequalities are strict, independently of label,
predicate and target. At arity 0 there are no inputs, so the allowed
premise set is empty. This gives the literal strict stratification
on the entire ground judgement set, not just on derived transitions.
The full ordinal/regular-cardinal result
Theorem~\ref{thm:p09-np-stratified} now applies: all rule families
are sets, all terms are finite, and premise cardinalities have a
common regular bound. It supplies modelhood, support, unique
stable and supported meaning and total well-foundedness. No
effective finite branching condition is needed for these conclusions.
\end{proof}
\end{annexitem}

\begin{annexitem}{annex:p09-sos-congruence}{thm:p09-sos-congruence}{Compatible closure is itself a bisimulation}
\begin{proof}
The union of bisimulations is a bisimulation: a pair comes from
one constituent and its witnesses remain in the union. Identity
is a bisimulation. The converse of one is a bisimulation by the
two-sided clauses. Composition is also a bisimulation: match a
move through the middle process in the first relation, then the
same-labelled middle move in the second; the two endpoint witnesses
compose. Predicate agreement composes as well, and reverse moves
are treated in reverse order. Thus the largest relation \(\sim\)
is reflexive, symmetric and transitive.

Let \(K_0=\sim\), and let \(K_{n+1}\) include \(K_n\)
and each \((f(\vec p),f(\vec q))\) whose argument pairs are
in \(K_n\). Put \(K=\bigcup_{n<\omega}K_n\). Finite arity
allows a common maximum of the finitely many argument stages,
so \(K\) is compatible. It is the least compatible relation
containing \(\sim\), by induction on stages. Structural induction
on any finite target term proves substitution transport: if every
variable's two images are \(K\)-related, the two instantiated
target terms are \(K\)-related. Variables are the base case;
constructors use compatibility. No transitivity of \(K\) is assumed.

Induct on the generation stage of a pair in \(K_n\) that its
predicates agree and every move on either side has a matching move
with endpoints in \(K\). Base pairs use \(\sim\). At a new
pair \((f(\vec p),f(\vec q))\), use support of the associated
model to select a ground root rule witnessing a move or predicate
on \(f(\vec p)\). Every positive action premise tests a child
\(p_i\). The induction hypothesis matches it from \(q_i\),
with endpoints in \(K\). Assign that matched endpoint to its
target variable. Distinct fresh target variables make all these
assignments consistent, even when a child has several action tests;
for an infinite premise set the stated Choice supplies the witness
family. Assign the input variables to \(q_i\).

If a negative action test held at \(p_i\), it holds at \(q_i\):
otherwise a \(q_i\)-move would have a reverse matching
\(p_i\)-move, contradiction. This argument excludes every
target, not only one chosen witness. Positive and negative
predicate tests transfer by predicate agreement. Thus all premises
of the same root rule hold under the new substitution. Modelhood
gives the matching head at \(f(\vec q)\). Its target is
\(K\)-related to the original target by the substitution lemma,
because the only variables are the inputs and assigned targets.
A predicate conclusion transfers directly. Reversing the argument
gives both directions. The induction proves that \(K\) itself
is a predicate-respecting bisimulation, hence \(K\subseteq\sim\).
Compatibility of \(K\) proves the desired congruence of \(\sim\).
The native nonrecursive GSOS language meets every stated condition;
its finiteness assumptions were not weakened by calling it another
format. Lookahead or recursive binders were not used in this proof.
\end{proof}
\end{annexitem}

\begin{annexitem}{annex:p09-sos-old-proof}{thm:p09-sos-old-proof}{Proof-height induction with the actual source-dependency distances}
\begin{proof}
In free first-order syntax, if \(\rho(t)\) is old, then every
variable occurring in \(t\) has an old image. A variable image
is a subterm of \(\rho(t)\); constructor recursion proves this
assertion. A fresh function name in \(t\) cannot disappear
under substitution. These facts use free syntax, not binding or
equational cancellation.

Induct on the well-founded height of a proof of a positive head
with old source, under negative hypotheses only. A positive head
cannot be a hypothesis leaf. If its root rule were new, a fresh
source could not instantiate to the old source. In the second
freshness case, the selected positive premise has an old source:
its variables occur in the root source and hence have old images,
and its function names are old. Its strictly smaller child proof
therefore, by induction, has old target and a label/predicate
occurring in \(T_0\). This contradicts the selected premise's
fresh target, action or predicate. Thus the root rule is old.

For that old rule, prove that every source-dependent variable has
old image. At dependency distance 0 this follows from the root
source. At a later distance it occurs in the target of a positive
premise whose source variables have already been assigned old
images. That premise source is old, and its smaller proof has
old target by the outer induction. Every variable image occurring
in that target is an old subterm, proving the next distance.
Every source-dependent variable has finite such distance: each
premise source term has finitely many variables, so the least
dependency closure is the union of its finite iterative stages.
This does not require a finite number of premises or a well-founded
variable dependency graph.

Source dependency now makes every occurring variable old. Hence
every used premise source/target and the head target are old.
All positive children have smaller old-source proofs, which transform
by outer induction to old proofs. Negative children are hypothesis
leaves from \(N\), and are old instances of the old rule's negatives.
The ground root instance is an old instance as well. If substitution
is taken to be total on the global variable supply, extend its
old assignments on the occurring variables by mapping unused
variables to one fixed old term; the assumed old root source
supplies such a term. This proves the entire old proof. If there
are no old closed terms, the assertion has no old-source case.
Both inductions are now justified and establish all stated claims.
\end{proof}
\end{annexitem}

\begin{annexitem}{annex:p09-sos-conservative}{thm:p09-sos-conservative}{Transfer of every negative-assumption rule, not only a selected relation}
\begin{proof}
Apply the preceding lemma to every proof
\(T_0\oplus T_1\vdash N/\alpha\) of the declared kind.
Its used negative leaves remain in \(N\), and its transformed
rule nodes are all ground old instances, so it is a proof from
\(T_0\) under those same hypotheses. Unused hypotheses do
not affect a proof tree. Conversely an old proof is an extended
proof, since signatures agree on common arities and the old rules
remain in the union. This gives both directions and exactly the
source's proof-level conservative-extension definition. The
counterexamples below verify why the two named hypotheses cannot
be omitted. No assumption about stable uniqueness was needed here.
\end{proof}
\end{annexitem}

\begin{annexitem}{annex:p09-sos-models}{thm:p09-sos-models}{Safe reducts, projection and an actual interval lift}
\begin{proof}
Call an extended relation safe if no action on an old source has
a new target. For every candidate \(M\), \(\Gamma_1(M)\)
is safe: every positive reduct proof restores to an ordinary proof
under negative hypotheses, and operational conservativity supplies
a proof from \(T_0\). A positive conclusion cannot be a negative
hypothesis leaf, so its root is an old ground rule, with old target
and old rule label/predicate. This uses proof-level conservativity
itself, not only the sufficient syntactic criterion.
For a safe candidate, let \(M_0\) be its old-source restriction.
Then
\(\Gamma_1(M)|_{\rm old}=\Gamma_0(M_0)\).
In the forward direction restore a reduct proof, use operational
conservativity, and note that each negative used in the old proof has old source
and old vocabulary. Its truth is the same in \(M\) and \(M_0\):
safety excludes all extra possible action targets. For the reverse,
an old reduct proof is an extended one with those same negative
truths. This proves equality, not an assumed monotonicity of negation.

Any three-valued stable pair \((L,U)\) has both components in
the image of \(\Gamma_1\), hence safe. The equality just proved
gives \(L_0=\Gamma_0(U_0),U_0=\Gamma_0(L_0)\), establishing
the projection claim. To lift a given old pair, start with lower
\(L_0\) and upper \(U_0\) together with all judgements on
new sources, excluding old-source/new-target atoms. Apply
\((L,U)\mapsto(\Gamma_1(U),\Gamma_1(L))\), with union
of lowers and intersection of uppers at limits. At the first
stage the lower contains \(L_0\), and the upper is contained
in the initial upper, by the reduct restriction equality. The
operator is monotone in the information order, so lowers increase
and uppers decrease. Consistency \(L\subseteq U\) persists
by antitonicity at successors; at a limit any lower member entered
earlier and remains in every later upper, hence in their intersection.
Safety and the exact fixed old lower/upper parts persist at all stages.

This iteration stabilizes: a strict information change classifies
an as-yet unclassified atom as true or false; no atom can subsequently
change that classification in an increasing consistent information
chain. A chain of length \(|B_1|^+\) of strict changes would
therefore inject that cardinal into the judgement set. At a stationary
stage the pair satisfies the two stable equations and has exactly
the desired old restriction. This constructs the lift without
claiming it unique. Finally iterate \(\Gamma_1^2\) from empty
to obtain the least lower. Every iterate is safe, restrictions
commute with successors by the equality above and with limit unions.
Its stationary old part is thus the least \(\Gamma_0^2\)
fixed point. Applying \(\Gamma\) gives the corresponding
old upper. This proves conservativity of the well-founded pair too.
\end{proof}
\end{annexitem}

\begin{annexitem}{annex:p09-sos-failures}{thm:p09-sos-failures}{Literal side-condition failures with distinguishing observations}
\begin{proof}
In the shared-target example, \(0\) and \(1\) are stopped.
The base constant \(p\) has both its \(a,b\)-moves to 0;
\(q\) has \(a\) to 0 and \(b\) to 1. Pair \(p,q\), their
converses and all pairs of stopped constants to obtain a bisimulation.
The shared variable \(y\) is instantiated to 0 at \(u(p)\),
giving its \(c\)-move. At \(u(q)\) the same variable would
have to be both 0 and 1, impossible in free syntax. No other rule
gives that move. Every other local-input requirement holds; only
target-variable distinctness fails. Thus this is a discriminating
case, not just a syntactically ill-formed rule list.

For the source pattern example, neither 0 nor \(0+0\) has a
move. Their two-way relation is a bisimulation. The exact axiom
for \(\zeta(0)\) gives its \(a\)-move, but cannot match
\(\zeta(0+0)\); there is no other \(\zeta\)-rule. This
positive TSS is stratified but not congruent. Thus stratification
alone cannot replace the format requirements.

For the actual source-dependency failure, the old signature contains
only constant \(a\), and the only old rule is \(xP/aP\).
Its ground instance requires \(aP\) to prove \(aP\), with
no well-founded proof. Adding the fresh constant axiom \(bP\)
allows the old rule with \(x=b\), proving \(aP\). The missing
variable \(x\) was not source-dependent. For the negative-only
hypothesis failure take old constant \(a\), no old rules, and
new constant \(b\). The registered new rule has old source
and a positive premise with fresh target \(b\), meeting the
second alternative. Under that premise as a positive hypothesis,
it proves an old-source conclusion, which the empty old TSS
cannot prove. Thus ordinary hypotheses cannot replace the theorem's
negative-only hypotheses.
\end{proof}
\end{annexitem}

\begin{annexitem}{annex:p09-sos-applications}{thm:p09-sos-applications}{Both freshness branches in the actual process applications}
\begin{proof}
In each BPA rule its input variables occur in its source. Each
derivative variable \(x'\) or \(y'\) occurs in a positive
action-premise target whose source is one of those inputs; it is
therefore acquired at the next dependency stage. There are no
other variables. This verifies all nine old rules, not only the
displayed left-sequence example.
Both priority rules have the fresh operator \(\theta\) in
their source, meeting the first alternative. In the exact timed
Table~3, the delay axiom has fresh \(\sigma_d\) in its source;
each of the other five rules has a positive timed premise on an
input variable in its source, with fresh action label \(\sigma\).
They meet the second alternative despite their old choice/product
heads. In the silent-step extension the atomic \(\tau\) axiom
has a fresh source constant. Its two new choice-action and two
new sequence-action rules each have a positive \(\tau\)-premise
on an input variable of the conclusion, meeting the second
alternative. Thus the complete conservative-extension theorem
applies to all three separate extensions.
The BPA, priority and timed rules also meet the declared local-input
format: input variables are distinct; derivative targets are fresh
and distinct; every action/predicate/negative premise tests an
input; no other variable is used. The earlier complete source
tables retain all success predicates and both versus one timed
branches. Therefore the proved congruence theorem applies to all
their operators. No recursive or later altered timed semantics
was included by this table-specific argument.
\end{proof}
\end{annexitem}

\begin{annexitem}{annex:p09-sos-panth-example}{thm:p09-sos-panth-example}{The exact source completeness counterexample}
\begin{proof}
The two axioms ground exactly \(aP,bP\). Other closed terms
cannot obtain \(P\), since no rule has such a head. Let
\(A=f(a)Q_1,B=f(a)Q_2,C=f(b)Q_1,D=f(b)Q_2\).
Only those four additional atoms can be grounded: every rule has
the positive test \(xP\), so its input must be \(a\) or \(b\).
After the two fixed axioms, the reduct generates precisely
\[
 B\iff\neg A\land\neg B,\quad
 D\iff\neg C\land\neg B,\quad
 A\iff\neg B\land\neg C,\quad
 C\iff\neg D\land\neg C.
\]
These are candidate-membership tests on the right. A stable fixed
point cannot contain \(B\) or \(C\), since each would require
its own absence. Their absence then forces \(A,D\), and those
values satisfy all four equations. This proves uniqueness of the
stated stable model on the full infinite term universe.
From empty the first reduct gives the two axioms and all four
\(Q\)-atoms; the next gives just the two axioms. Squared
iteration thereafter is stationary at those axioms, with all four
\(Q\)-atoms in the upper bound. Thus they are unknown and
the TSS is not complete. It meets panth: its only premises of
action type are absent, its source \(f(x)\) has one function
symbol and a single source variable. Predicate and negative
premises are allowed. In the unique stable model, \(a,b\)
have the same predicates and no moves, so are bisimilar, while
their \(f\)-images disagree on \(Q_1,Q_2\). This is the
source's exact missing-completeness failure, not a counterexample
to the complete-panth theorem or to the local-input theorem.
\end{proof}
\end{annexitem}
